\documentclass[11pt,a4paper]{../_shared/altacv}

\geometry{left=1.5cm,right=1.5cm,top=1.cm,bottom=1.2cm,columnsep=1.5cm}

\usepackage[utf8]{inputenc}
\usepackage[T1]{fontenc}
\usepackage[french]{babel}
\usepackage{paracol}
\usepackage{xcolor}
\usepackage{hyperref}
\usepackage{fontawesome5}
\usepackage{setspace}

\definecolor{SlateGrey}{HTML}{2E2E2E}
\definecolor{LightGrey}{HTML}{666666}
\definecolor{DarkPastelRed}{HTML}{8AC0D9}
\definecolor{PastelRed}{HTML}{00B0DA}
\definecolor{GoldenEarth}{HTML}{B4AD0C}
\colorlet{name}{black}
\colorlet{tagline}{PastelRed}
\colorlet{heading}{DarkPastelRed}
\colorlet{headingrule}{GoldenEarth}
\colorlet{subheading}{PastelRed}
\colorlet{accent}{PastelRed}
\colorlet{emphasis}{SlateGrey}
\colorlet{body}{LightGrey}

\renewcommand{\familydefault}{\sfdefault}

\name{\Large Guillaume BOILEAU}
\tagline{Chef de projet technique | Systèmes complexes, logiciel, data, validation et coordination}
\personalinfo{
  \email{guillaume.boileaupro@gmail.com}
  \phone{+33 6 59 94 85 84}
  \location{Alpes-Maritimes, France}
  \linkedin{guillaume-boileau-phd-891b81265}
  \researchgate{Guillaume-Boileau-2}
  \orcid{0000-0002-3576-6968}
}

\setlength{\parskip}{0.75\baselineskip}

\begin{document}
\makecvheader
\vspace{-0.6cm}

\cvsection{Lettre de motivation}
\vspace{-0.4cm}

{\color{accent}\textbf{Objet :}} \textit{Candidature spontanée - Chef de projet technique / Ingénieur systèmes, logiciel, data et validation}

{\color{body}

Madame, Monsieur,

Je vous adresse ma candidature spontanée pour un poste de chef de projet technique, ingénieur systèmes, ingénieur logiciel, ingénieur IVVQ ou coordinateur technique sur des projets à forte composante logicielle, data, spatiale ou industrielle.

Mon parcours s'est construit à l'interface entre systèmes complexes, développement logiciel, analyse de données, validation, coordination technique et pilotage de contributions. J'ai travaillé dans des environnements exigeants où la robustesse des outils, la qualité des livrables, la traçabilité, la structuration des besoins et la coordination entre plusieurs acteurs sont essentielles.

Au Laboratoire Lagrange de l'Observatoire de la Côte d'Azur, dans le cadre de travaux menés avec le CNES, j'ai contribué aux activités de développement, d'intégration et de validation pour la mission spatiale LISA, mission ESA/NASA de grande envergure. J'y ai notamment développé CATpyLISA, une plateforme Python destinée à l'intégration, la comparaison et la validation de modèles et de chaînes d'analyse complexes. Ce travail m'a amené à structurer des workflows techniques, à organiser des campagnes de comparaison, à analyser des incohérences, à documenter les résultats et à coordonner des contributions dans un cadre international.

Cette expérience m'a permis de développer une posture directement transposable à des fonctions de chef de projet technique : comprendre rapidement un système, clarifier les objectifs, structurer les actions, organiser les priorités, suivre les contributions, sécuriser les livrables et maintenir une vision d'ensemble sur des sujets complexes. J'ai également encadré plusieurs stagiaires de niveau Master et coordonné des activités impliquant ingénieurs, techniciens et chercheurs, avec un souci constant de lisibilité, de rigueur et de continuité dans l'exécution.

Plus récemment, chez VEV Platform Services France, j'ai travaillé sur des services logiciels liés aux infrastructures de recharge électrique, avec des interactions entre plateforme logicielle, flux de données opérationnels et systèmes terrain. Cette expérience industrielle m'a permis de renforcer ma capacité à travailler sur des services en conditions réelles, à analyser des incidents, à contribuer à la fiabilité opérationnelle, à suivre des échanges de données et à intervenir dans une logique d'amélioration continue.

Je souhaite aujourd'hui mettre cette double culture, issue du spatial et du logiciel industriel, au service de projets techniques ambitieux. Je peux intervenir sur des sujets de coordination, pilotage d'activités, suivi d'exigences, validation, automatisation, analyse d'anomalies, documentation technique, structuration de workflows et interface entre équipes métier, logiciel, système et management.

Je suis particulièrement à l'aise dans les environnements où il faut faire le lien entre des problématiques techniques complexes et des objectifs opérationnels concrets. Mon expérience m'a appris à transformer des besoins parfois hétérogènes en actions structurées, à produire des éléments de décision clairs, à fiabiliser des processus et à accompagner des équipes vers un objectif commun.

Rigoureux, autonome et habitué aux environnements internationaux, je suis motivé par les postes où la valeur ajoutée repose autant sur la compréhension technique que sur la capacité à coordonner, prioriser, documenter et faire avancer les sujets. Je souhaite m'investir dans une structure où je pourrai contribuer à des projets exigeants, utiles et techniquement solides.

Je serais heureux de pouvoir échanger avec vous afin de vous présenter plus en détail mon parcours et la manière dont je pourrais contribuer efficacement à vos projets.

Je vous prie d'agréer, Madame, Monsieur, l'expression de mes salutations distinguées.

\textbf{Guillaume Boileau}

}

\end{document}